\documentclass{article}
\usepackage[utf8]{inputenc}
\usepackage[T1]{fontenc}
\usepackage{geometry}
\usepackage{graphicx}
\usepackage{float}
\usepackage{hyperref}
\usepackage[english,indonesian]{babel}

\title{Optimasi Pemrosesan Template LaTeX Menggunakan Pendekatan Berbasis AI}
\author{Academic Support Agent}
\date{\today}

\begin{document}
\maketitle

\begin{abstract}

\end{abstract}

\chapter{PENDAHULUAN}


\chapter{TINJAUAN PUSTAKA}


% Original content (commented out):
% % Thesis Template HRW
% 2022-11-25 v. 0.1.0
% Lucas Gerritzen
% Hochschule Ruhr West
\input{Einstellungen}

\begin{document}
	\include{Kapitel/Titel}
	
	\frontmatter
	\pagenumbering{Roman}
	
	\include{Kapitel/Kurzfassung}
	
	\include{Kapitel/Sperrvermerk}
	
	\include{Kapitel/Erklaerung}

	\tableofcontents
	\addcontentsline{toc}{chapter}{Inhaltsverzeichnis}
	\listoffigures
	\addcontentsline{toc}{chapter}{Abbildungsverzeichnis}
	\listoftables
	\addcontentsline{toc}{chapter}{Tabellenverzeichnis}
	
    \include{Verzeichnisse/Abkürzungsverzeichnis}

	\mainmatter
	\lhead{}
	\chead{}
	\pagestyle{fancy}

	\include{Kapitel/Einleitung}

	%\include{chapter2}

	%\include{chapter3}
	
	\include{Kapitel/Zusammenfassung_und_Ausblick}
	
	\include{Kapitel/Anhang}
	
	\printbibliography
	\addcontentsline{toc}{chapter}{Literaturverzeichnis}

\end{document}


\end{document}